% Generated draft: requires mathematical and source review.
\documentclass[11pt]{article}
\usepackage[margin=25mm]{geometry}
\usepackage{fontspec}
\setmainfont{Latin Modern Roman}
\usepackage{amsmath,amssymb,mathtools,unicode-math}
\setmathfont{Latin Modern Math}
\usepackage{enumitem,needspace,xurl,hyperref}
\setlength{\parindent}{0pt}
\setlength{\parskip}{5pt}
\setlength{\emergencystretch}{3em}
\newcommand{\N}{\mathbb{N}}
\newcommand{\Z}{\mathbb{Z}}
\newcommand{\Q}{\mathbb{Q}}
\newcommand{\R}{\mathbb{R}}
\newcommand{\C}{\mathbb{C}}
\newcommand{\F}{\mathbb{F}}
\newcommand{\A}{\mathbb{A}}
\newcommand{\PP}{\mathbb{P}}
\newcommand{\E}{\mathbb{E}}
\newcommand{\eps}{\varepsilon}
\newcommand{\dd}{\,\mathrm{d}}
\begin{document}
\section*{UnsolvedMath ID: 30000649}
\textbf{Source problem code:} catalogue entry
\subsection{Definitions and mathematical statement}
For \(n\in\Z_{>0}\), the complex Weyl algebra is the unital associative algebra
\[
 A_n(\C)=\C\langle X_1,\ldots,X_n,D_1,\ldots,D_n\rangle
 \big/\bigl([X_i,X_j],\ [D_i,D_j],\ [D_i,X_j]-\delta_{ij}1\bigr)_{1\leq i,j\leq n},
\]
where \([a,b]=ab-ba\) and \(\delta_{ij}\) is the Kronecker symbol. The angle brackets denote the free associative algebra. This makes explicit that the source's polynomial-style notation is subject to noncommutative Weyl relations: the \(X_i\) commute among themselves and the \(D_i\) commute among themselves, but the two families do not all commute.

The Dixmier conjecture, named in the original record, asks whether
\[
 \forall n\geq1\ \forall\phi\in\operatorname{End}_{\C\text{-}\mathrm{alg}}(A_n(\C)),
 \qquad \phi\in\operatorname{Aut}_{\C\text{-}\mathrm{alg}}(A_n(\C)).
\]
Endomorphisms are unital \(\C\)-algebra maps and fix the scalar field pointwise. The conclusion asks for a \(\C\)-algebra map \(\psi\) with \(\psi\phi=\phi\psi=\operatorname{id}\). No assumption of degree preservation, continuity, filtration compatibility or prior surjectivity is made.

Equivalently, images \(Q_i=\phi(X_i)\), \(P_i=\phi(D_i)\) must satisfy
\[
 [Q_i,Q_j]=[P_i,P_j]=0,\qquad [P_i,Q_j]=\delta_{ij}1,
\]
and the question is whether every resulting unital endomorphism has an inverse. These are elements of the Weyl algebra, not merely commuting polynomial functions with a prescribed Jacobian. The field and the number of generator pairs remain those in the source; no passage to another characteristic, completion or additional generators is substituted.

The literature triage checked 2026-08-21 reports no authoritative complete resolution of the extracted conjecture located in its search. The open label and verification note reporting removal of the label and notation cleanup remain dated metadata, not independently verified current status.
\subsection{Short English statement}
Does every unital complex-algebra endomorphism of each Weyl algebra have an inverse? The noncommutative relations, fixed scalar field and unrestricted endomorphism class are explicit. No filtration-preserving or commutative-polynomial version is substituted.
\subsection{Sources}
\begin{itemize}
\item Local UnsolvedMath registry, permanent ID 30000649, source code catalogue entry, ``Dixmier Conjecture for Weyl-Algebra Endomorphisms.'' The original statement names the Dixmier conjecture.
\item \emph{Affine Algebraic Geometry (2007)}, Oberwolfach Reports, pp.\ 5--82, DOI 10.4171/owr/2007/01: \url{https://doi.org/10.4171/owr/2007/01}. This is also the sole supplied literature reference.
\item Literature triage checked 2026-08-21; local record updated 2026-09-16. No external source has been visited or independently checked here.
\end{itemize}
\end{document}
